


\Large{Ilova 1. Xalqaro sistema (SI) ning asosiy va qo'shimcha birliklari}
\begin{table}[!ht]
\centering
\begin{tabular}{lll}
\toprule
\textbf{Kattalik} & \textbf{Birlik nomi} & \textbf{Belgisi} \\
\midrule
\multicolumn{3}{l}{\textbf{Asosiy birliklar}} \\
\midrule
Uzunlik & metr & m \\
Massa & kilogramm & kg \\
Vaqt & sekund & s \\
Elektr tokining kuchi & amper & A \\
Termodinamik harorat (temperatura) & kelvin & K \\
Modda miqdori & mol & mol \\
Yorug'lik kuchi & kandela & kd \\
\midrule
\multicolumn{3}{l}{\textbf{Qo'shimcha birliklar}} \\
\midrule
Yassi burchak & radian & rad \\
Fazoviy burchak & steradian & sr \\
\bottomrule
\end{tabular}
\end{table}





\Large{Ilova 2. Quyosh, Yer va Oyga taalluqli ba'zi o'lchamlar}
\begin{table}[!ht]
\centering
\begin{tabular}{llll}
\toprule
\textbf{O'lchamlar} & \textbf{Quyosh} & \textbf{Yer} & \textbf{Oy} \\
\midrule
Massa, kg & $1,97 \cdot 10^{30}$ & $5,96 \cdot 10^{24}$ & $7,33 \cdot 10^{22}$ \\
Radius, m & $6,95 \cdot 10^8$ & $6,37 \cdot 10^6$ & $1,74 \cdot 10^6$ \\
O'rtacha zichligi, $\text{kg/m}^3$ & 1400 & 5518 & 3350 \\
Yergacha o'rtacha masofa, km & $1,496 \cdot 10^8$ & --- & 384440 \\
\bottomrule
\end{tabular}
\end{table}







\Large{Ilova 3. Gazlar, suyuqliklar va qattiq jismlarning zichliklari ($\text{kg/m}^3$)}

\begin{table}[!ht]
\centering
\small
\begin{tabularx}{\textwidth}{Xl | Xl | Xl}
\toprule
\multicolumn{2}{c|}{\textbf{Gazlar (0 $^\circ\text{C}$ da)}} & \multicolumn{2}{c|}{\textbf{Suyuqliklar}} & \multicolumn{2}{c}{\textbf{Qattiq jismlar}} \\
\midrule
\textbf{Modda} & \textbf{Zichlik} & \textbf{Modda} & \textbf{Zichlik} & \textbf{Modda} & \textbf{Zichlik} \\
\midrule
Vodorod & 0,08988 & Spirt & 790 & Po'kak & 200 \\
Havo & 1,293 & Kerosin & 800 & Muz & 900 \\
Kislorod & 1,429 & Benzol & 880 & Osh tuzi & 2200 \\
Karbonat angidrid & 1,977 & Kanakunjut moyi & 950 & Alyuminiy & 2600 \\
& & Suv ($+4 ^\circ\text{C}$) & 1000 & Shisha & 2700 \\
& & Qon & 1050 & Kaliy & 7100 \\
& & Glitserin & 1200 & Po'lat & 7700 \\
& & Kumush (suyuq) & 9300 & Temir & 7900 \\
& & Simob & 13600 & Jez & 8400 \\
& & Oltin (suyuq) & 17200 & Mis & 8600 \\
& & & & Nikel & 8800 \\
& & & & Kumush & 10500 \\
& & & & Qo'rg'oshin & 11300 \\
& & & & Oltin & 19300 \\
& & & & Platina & 21400 \\
\bottomrule
\end{tabularx}
\caption*{Eslatma: Gazlar zichligi normal atmosfera bosimi ($1,013 \cdot 10^5$ Pa) sharoitida berilgan.}
\end{table}










\Large{Jadval 4. Suvning turli haroratlardagi zichligi}

\begin{table}[!ht]
\centering
\begin{tabular}{cc}
\toprule
\textbf{Harorat, $^\circ\text{C}$} & \textbf{Zichlik, $\text{kg/m}^3$} \\
\midrule
20 & 998 \\
30 & 996 \\
40 & 992 \\
50 & 988 \\
60 & 983 \\
70 & 978 \\
80 & 972 \\
\bottomrule
\end{tabular}
\end{table}
